% page 91 of the facsimile (image p-090 of the scan)
% block 182.9 x 237.7 mm ; left margin 19.1 mm
\pgc{-2.540mm}{0.000mm}{
\l{\cel{59}83}
\l{}
\l{\cel{6}\sou{centaureo}. (\sou{bot}.)Planto de la familio \textquotedbl{}kompozaji\textquotedbl{}, di qua plura}
\l{\cel{9}speci uzesas medicinale. - L.\cel{1}\sou{chlora perfoliata centaurea}.}
\l{\cel{9}- EFIS.}
\l{}
\l{\cel{6}\sou{centezimala}. Di qua la parti esas centimi. - DEFIRS.}
\l{}
\l{\cel{6}\sou{centiaro}. La parto centima di aro. - DEFIRS.}
\l{}
\l{\cel{6}\sou{centifolio}. (\sou{bot.}) Rozo di qua la petali esas tre multa.}
\l{\cel{9}L.\cel{1}\sou{rosa centifolia}. - DEF.}
\l{}
\l{\cel{6}\sou{centigrado}.\cel{2}Dividita a 100 gradi. - DEFIS.}
\l{}
\l{\cel{6}\sou{centigramo}. Parto centima di gramo. - DEFIRS.}
\l{}
\l{\cel{6}\sou{centilitro}. Parto centima di litro. - DEFIRS.}
\l{}
\l{\cel{6}\sou{centimo.}\cel{1}Parto centima di franko (la unajo monetala di Francia).}
\l{\cel{9}- DEFIRS.}
\l{}
\l{\cel{6}\sou{centimetro.}\cel{2}Parto centima di metro. - DEFIRS.}
\l{}
\l{\cel{6}\sou{centono.}\cel{2}Versajo ek versi o fragmenti di versi pruntita hike}
\l{\cel{9}ed ibe. - DEF.}
\l{}
\l{\cel{6}\sou{centro}. I. (\sou{geom}.) Punto interna, situita egaldiste de omna}
\l{\cel{9}punti di cirklokurvo o di la surfaco di sfero. - II.To quo}
\l{\cel{9}esas la centro, la mezo di extensajo. - II. Apliko-punto di}
\l{\cel{9}la rezultanto de plura forci. - IV. Loko adube konvergas forci,}
\l{\cel{9}agi diversa. - DEFIRS.}
\l{}
\l{\cel{6}\sou{centrifugar}. (\sou{trans.}) Efikigar la forco per qua moveblo qua}
\l{\cel{9}cirkulas segun kurvo retropulsesas de la centro di ta kurvo.-}
\l{\cel{9}DEFIRS.}
\l{}
\l{\cel{6}\sou{centripetar}. (\sou{trans.}) Efikigar la forco per qua moveblo qua}
\l{\cel{9}cirkulas segun kurvo atraktesas ad la centro di ta kurvo. -}
\l{\cel{9}DEFIRS.}
\l{}
\l{\cel{6}\sou{centrolecito}.\cel{2}Ovo di qua la plasmo nutriva akumulesas en la}
\l{\cel{9}centro di la celulo. - DEFIRS.}
\l{}
\l{\cel{6}\sou{centrosfero}. (\sou{biol.}) Maso de protoplasmo qua cirkondas la centro-}
\l{\cel{9}somo e prizentas ulakaze du zoni sama-centra e neagale kolo-}
\l{\cel{9}roza. - DEFIRS.}
\l{}
\l{\cel{6}\sou{centrosomo}. (\sou{biol.}) Mikra korpuskulo qua okupas ordinare la}
\l{\cel{9}centro di la celulo, apud la parieto nukleala, e qua konjekt-}
\l{\cel{9}eble pleas granda rolo en mitoso. - DEFIRS.}
\l{}
\l{\cel{6}\sou{centurio}. (\sou{En Roma antiqua}): I. Asemblitaro de cent civitani ye}
\l{\cel{9}qui konsistis un ek la dividuri politikala di la populo.}
\l{\cel{9}- II. Dividuro de kohorto a cent homi. - DEFIRS.}
}
